% !TeX program = latexmk
\documentclass[UTF8,a4paper,zihao=5,fontset=none]{ctexart}
\usepackage[a4paper,margin=2.25cm,headheight=15pt]{geometry}
\setCJKmainfont{FandolSong-Regular.otf}[cmap=UniGB-UTF16-H,BoldFont=FandolHei-Regular.otf,ItalicFont=FandolKai-Regular.otf]
\setCJKsansfont{FandolHei-Regular.otf}[cmap=UniGB-UTF16-H]
\setCJKmonofont{FandolFang-Regular.otf}[cmap=UniGB-UTF16-H]
\usepackage{float}
\usepackage{amsmath,amssymb,booktabs,tabularx,graphicx,fancyhdr,pgfplots}
\usepackage[hidelinks,unicode]{hyperref}
\pgfplotsset{compat=1.18}
\setlength{\parskip}{.3em}\linespread{1.12}
\setlength{\emergencystretch}{2em}
\pagestyle{fancy}\fancyhf{}\fancyhead[L]{\small B 题问题二：期望包围圆面积模型}\fancyhead[R]{\small R001 数值求解}\fancyfoot[C]{\thepage}
\newcommand{\E}{\mathbb E}\newcommand{\Prob}{\mathbb P}
\title{第二检测点的数值求解报告\\\large 原点、正东基准情形；含候选区域与独立核验}
\author{}\date{2026 年 9 月 11 日}
\begin{document}
\maketitle\vspace{-2em}
\section{结论与适用条件}
按照已保存的《第二问模型》，将第二次无线电检测结束时、源位置后验支持集的最小包围圆面积的期望作为优化目标。第一次检测点取 \(S_1=(0,0)\)，已测得正常示向度 \(\theta_1=0^\circ\)。经过全域网格搜索及逐级加密，得到数值最优点约为
\[
 q^*=(932.6,\ 590.0)\ \mathrm m
 \quad\text{或}\quad (932.6,\ -590.0)\ \mathrm m.
\]
两个点关于正东方向对称，基线长度约为 \(1103.56\) 米，方向为正东向北或向南偏转 \(32.319^\circ\)。以下所有主模型指标按未取整点 \((932.6,590)\) 计算；坐标写作 \((933,590)\) 时只是行动层面的取整建议。

\begin{table}[H]\centering\small
\caption{两个优化问题的数值结果；面积单位为平方米，半径单位为米}
\begin{tabular}{lrr}\toprule
指标 & 允许全部反馈的主模型 & 附加保证接收约束 \\\midrule
第二点横坐标 & 932.600 & 864.316 \\
第二点纵坐标（北侧） & 590.000 & 511.356 \\
期望最小包围圆面积 \(J\) & 2103.7834 & 2228.8321 \\
包围圆半径的期望 \(\E[\rho]\) & 25.1280 & 25.3032 \\
无信号概率 \(P_N\) & 0.067262\% & 0 \\
达到 \(\rho\le20\) 的概率 \(P_{20}\) & 20.6622\% & 34.1389\% \\
强信号概率 \(P_H\) & 0 & 0 \\\bottomrule
\end{tabular}
\end{table}
主模型允许极小概率的无信号反馈，以换取较低的平均面积。若要求第二次一定能够接收，最优期望面积增加约 \(5.9440\%\)，但满足 \(20\) 米保证定位条件的概率反而上升。这说明\textbf{最小化平均包围圆面积与最大化一次后续清除的保证成功概率是不同目标}。这里的 \(P_{20}\) 指第二次反馈后存在一个圆心，从该处到全部可行源位置都不超过 \(20\) 米的概率；它既不是在第二检测点立即清除的概率，也不是整体任务的清除率。

\subsection{题目事实与额外假设}
原题要求为一个全向源选择第二检测点并给出候选区域。目标分布圆域半径为 \(1800\) 米，角误差界为 \(\pm1^\circ\)，接收半径为 \(1000\) 至 \(1500\) 米，同一地点的误差固定，\(5\) 米内为强信号，\(20\) 米内可光学定位，移动速度为 \(5\) 米每秒。第二问本身没有指定位置、角误差或半径的概率密度。

本次沿用已确认的建模假设：源位置按面积均匀；不同检测点的角误差在给定源与半径后独立，且均匀分布于 \([-1^\circ,1^\circ]\)；同一源的固定接收半径先验为 \(U[1000,1500]\)，与位置先验独立。使用平面直线移动，不附加障碍物；最小基线取 \(1\) 米，仅用于排除重复检测位置。半径先验和空间独立性是建模约定，不能表述为原题已知的随机生成规则。

\newpage
\section{计算对象与数值方法}
\subsection{第一次观测后的正确分布}
设 \(r=\|g\|\)、\(\alpha=\pi/180\)，令 \(w(r)=\Prob(R\ge r)\)。在 \(r\le1000\) 时 \(w(r)=1\)，在 \(1000<r<1500\) 时 \(w(r)=(1500-r)/500\)。第一次正常测向后的面积密度及其极坐标形式为
\[
 p_1(g)=\frac{w(\|g\|)}{Z_1}\mathbf1_{K_1}(g),\qquad
 p_1(r,\varphi)=\frac{r\,w(r)}{Z_1},\quad
 5<r<1500,\quad |\varphi|\le\alpha.
\]
这里 \(K_1\) 是上述扇环的闭包，圆域半径 \(1800\) 在本基准中不产生额外裁切；\(Z_1=\alpha\{\E[R^2]-25\}=27633.9434909\)。因正常接收事件的筛选作用，\(1000\) 米以外的位置密度下降，源距也不服从均匀分布。给定源位置后，同一个半径的后验为 \(R\mid g,\mathcal I_1\sim U[\max(1000,r),1500]\)。第二次检测必须沿用这个半径，而不是独立重新抽取。

\subsection{三类反馈及目标积分}
设 \(d=\|g-q\|\)，记 \(L_z(q)=\pi\rho^2(K_z(q))\)。强信号、无信号及正常示向度反馈分别记作 \(H,N,\theta\)，积分中对应的未归一化位置权重为
\[
\begin{aligned}
 f_H(g;q)&=\mathbf1_{K_1}(g)\mathbf1_{d\le5}\,w(r),\\
 f_N(g;q)&=\mathbf1_{K_1}(g)\{w(r)-w(\max(r,d))\},\\
 f_\theta(g;q)&=\frac{\mathbf1_{K_1}(g)\mathbf1_{d>5}}{2\alpha}
 \mathbf1_{|\operatorname{wrap}(\theta-\arg(g-q))|\le\alpha}\,w(\max(r,d)).
\end{aligned}
\]
对位置积分后除以 \(Z_1\)，分别得到 \(P_H,P_N\) 与示向度预测密度 \(p_\Theta\)。每个反馈的支持集取对应权重正值区域的闭包，几何边界保留。最终计算
\[
 J(q)=P_H(q)L_H(q)+P_N(q)L_N(q)
       +\int_0^{2\pi}p_\Theta(\theta\mid q,\mathcal I_1)L_\theta(q)\,\mathrm d\theta.
\]
在主模型最优点，无信号分支的包围圆半径约为 \(60.9207\) 米，对期望面积贡献 \(7.8424\) 平方米。强信号分支的概率为零；程序仍完整实现此分支，并在存在强信号的对照点进行了检查。

\subsection{几何与积分实现}
源支持集由圆盘、圆盘外部及半平面相交构成。程序求各边界交点，将圆周分为圆弧、直线分为线段，再用区间中点判断可行性。无信号分支采用 \(d\ge1000\) 与 \(d\ge r\)，后者化为半平面 \(2q\cdot g\le\|q\|^2\)。正常分支同时保留 \(5\) 米排除区与 \(1500\) 米接收上限。

对这些连续边界，先计算有限点集的最小包围圆，再在每条圆弧上精确搜索离圆心最远的位置，持续加入违反覆盖条件的点。有限子集给出半径下界，整个边界的最远距离给出覆盖半径上界，二者达到容差才停止。因此没有用稀疏源点云代替完整支持集。

位置积分采用极坐标射线分段及 Gauss--Legendre 积分，并在接收权重变化处切分；反馈角度积分按几何事件切分后自适应加密。计算 \(P_{20}\) 时额外求解 \(\rho=20\) 的角度阈值，避免用一般平滑函数的积分设置跨越指示函数跳点。

\newpage
\section{候选区域、搜索过程与策略解释}
令 \(\mathcal C\) 为距 \(K_1\) 不超过 \(1500\) 米且基线不小于 \(1\) 米的点集。离 \(K_1\) 更远的点必定无信号，也不能进一步缩小支持集，因而不优于已找到的点。利用关于横轴的对称性，仅搜索 \(q_y\ge0\)，再作镜像。

搜索先使用全域 \(100\) 米网格，再使用 \(50\) 米网格，共评估 \(2381\) 个有效细网格点；随后在最优区域内按 \(10,2,0.5,0.1\) 米加密。最终 \(0.1\) 米网格含 \(121\) 个点，采用高精度积分，并对选中的点作更高阶复算。保证接收的模型另外扫描可行域边界并检查内部点；北侧最优边界可写为 \(q=(5\cos\alpha,-5\sin\alpha)+1000(\cos t,\sin t)\)，在约 \(t=30.76^\circ\) 处达到最小值。

以相对最优值允许的损失定义候选区域 \(\mathcal C_\eta=\{q\in\mathcal C:J(q)\le(1+\eta)J(q^*)\}\)。图中给出北侧候选点，南侧取其镜像。
\begin{figure}[H]\centering
\begin{tikzpicture}
\begin{axis}[width=13.5cm,height=8.2cm,xmin=820,xmax=1030,ymin=450,ymax=735,xlabel={东向坐标 \(q_x\)（米）},ylabel={北向坐标 \(q_y\)（米）},grid=major,grid style={gray!18},tick label style={font=\small},legend style={at={(.99,.99)},anchor=north east,font=\small,fill=white,draw=gray!30},axis on top]
\addplot[only marks,mark=square*,mark size=1.6pt,color=blue!22] table[x=x,y=y,col sep=comma]{results/candidates_5pct.csv};
\addlegendentry{面积损失不超过最优值的 105\%}
\addplot[only marks,mark=square*,mark size=1.6pt,color=teal!70] table[x=x,y=y,col sep=comma]{results/candidates_1pct.csv};
\addlegendentry{面积损失不超过最优值的 101\%}
\addplot[only marks,mark=star,mark size=3.5pt,color=red!75!black]coordinates{(932.6,590)};
\addlegendentry{主模型数值最优点}
\addplot[only marks,mark=diamond*,mark size=3pt,color=orange!90!black]coordinates{(864.3164,511.3558)};
\addlegendentry{保证接收的数值最优点}
\end{axis}
\end{tikzpicture}

\caption{候选区域的 \(5\) 米网格近似。深色对应 \(\eta=1\%\)，浅色对应 \(\eta=5\%\)。绘制的是通过目标值检验的网格点，连续边界由相应等值线定义。}
\end{figure}

\(1\%\) 候选集在北侧含 \(267\) 个网格点，其坐标投影为 \(895\le q_x\le970\)、\(535\le q_y\le645\) 米；\(5\%\) 候选集含 \(1285\) 个点，坐标投影为 \(840\le q_x\le1010\)、\(465\le q_y\le715\) 米。\textbf{投影范围不是可任意组合的矩形候选区}，应使用图中的点集或附带 CSV，并按目标函数复核网格之间的位置。

直接沿示向度向东走到 \((1000,0)\) 时，期望面积约为 \(332380.28\) 平方米；选 \((750,750)\) 时约为 \(2463.45\) 平方米。适度横向偏移能改善两次方向约束的交角，而过大的横向移动又会增加远距离或接收失败风险。最优位置是后验源距权重、交会几何和反馈分支代价共同作用的结果，不能单由一个固定交角决定。

\newpage
\section{数值核验与敏感性}
\subsection{收敛及独立检查}
表中所有精度级别都在同一个点 \((932.6,590)\) 上计算，避免把优化点变化与积分误差混在一起。级别 3 开始对 \(P_{20}\) 的阈值单独切分；级别 2 的该项仅用于粗筛，不能作为最终结果。
\begin{table}[H]\centering\small
\begin{tabular}{crrr}\toprule
精度级别 & \(J\)（平方米）& \(P_{20}\) & 总概率与 1 的差的绝对值 \\\midrule
2 & 2103.687635 & 21.442889\% & \(2.29\times10^{-6}\) \\
3 & 2103.783444 & 20.662192\% & \(2.64\times10^{-10}\) \\
4 & 2103.783406 & 20.662192\% & \(1.21\times10^{-11}\) \\
5 & 2103.783403 & 20.662192\% & \(9.06\times10^{-13}\) \\\bottomrule
\end{tabular}
\end{table}
最后两级期望面积相差 \(2.81\times10^{-6}\) 平方米；镜像点的面积差小于 \(4\times10^{-11}\) 平方米。第一次扇环的归一化常数与最小包围圆还可解析检查，后者半径为 \(747.615393\) 米。

几何独立核验使用另一种表示：固定第一扇形内的角度，从原点发射射线，直接求各条件在半径方向上的可行区间。每个案例取 \(12001\) 条射线，从独立生成的端点计算覆盖残差与直径下界。共检查 \(47\) 个案例，其中 \(46\) 个有非空支持集，包含正常、强信号、无信号及保证接收的退化无信号案例。最大采样覆盖残差约 \(6.60\times10^{-12}\) 米，主算法半径与独立直径下界的最大差为 \(0.002341\) 米。此检查为离散独立核验，不替代对全部连续输入的形式化证明。

概率核验直接从 \(p_1\) 采样源位置，再从条件截断分布采样同一个 \(R\)，最后生成第二次反馈。在最优点的 \(500000\) 次独立模拟中，平均面积为 \(2099.6735\) 平方米，标准误为 \(1.6435\) 平方米，与积分结果相差 \(2.50\) 个标准误；\(P_{20}\) 的模拟值为 \(20.8024\%\)，与积分结果相差约 \(2.45\) 个标准误。保证接收点另作 \(500000\) 次模拟，平均面积为 \(2231.3114\)，标准误为 \(2.3425\)，偏差为 \(1.06\) 个标准误。全部实际采样源都落在对应的可行支持集及覆盖圆内。概率采样验证和几何射线验证各自检查不同环节，不能互相替代。

\subsection{接收半径先验敏感性}
保持 \(1000\) 至 \(1500\) 米的支持不变，分别用密度随半径线性下降及线性上升的三角先验重新优化。它们是对不确定先验的敏感性检查，不是新增题目事实。
\begin{table}[H]\centering\small
\begin{tabular}{lrrr}\toprule
半径先验 & 北侧数值最优点（米） & 最优期望面积 & 沿用基准点的相对损失 \\\midrule
偏向小半径 & \((878,570)\) & 1801.8711 & 2.48\% \\
均匀 & \((932.6,590)\) & 2103.7834 & 0 \\
偏向大半径 & \((980,608)\) & 2265.6626 & 1.57\% \\\bottomrule
\end{tabular}
\end{table}
偏向小半径时应适当收短基线，偏向大半径时可向东移动。这支持把本结论写成\textbf{明确先验下的数值策略}。此外，\(P_{20}\) 对位置与先验也较敏感，最小面积策略不能自动成为第三问的最短时间策略。

\newpage
\section{交付内容、复现方式与结论边界}
\subsection{文件组织}
本次计算全部存放在新的 \path{R001} 文件夹中，用户的模型文档保持原样。\path{model_snapshot.tex} 是本次采用模型的快照，\path{config.json} 记录输入与模型散列值，\path{src/solver.cpp} 是不依赖外部数值库的 C++17 求解器。\path{results/summary.json} 是最终机器可读汇总；\path{comparison.csv}、\path{candidates_1pct.csv}、\path{candidates_5pct.csv} 和各级网格结果保存对应数值。\path{verification} 保存收敛、蒙特卡洛、独立几何和语义检查记录。

\path{src/reproduce.py} 可重算主要结果和验证，完整网格复现使用 \path{--full} 参数；射线几何核验另依赖 Python 的 NumPy。已保存所有网格输入，随机模拟的样本数和种子写入结果文件。运行方式详见同目录的 \path{README.md}。

\subsection{如何使用本结果}
如果严格沿用当前模型，取 \((932.6,\pm590)\) 作为第二检测点即可。在北侧和南侧没有额外环境差异时，两者等价。若要缩短移动时间，可在 \(1\%\) 或 \(5\%\) 候选集中按基线长度进行第二级排序，而不是另设一个没有量纲依据的面积与时间加权和。主最优点的移动时间约为 \(220.71\) 秒，之后正常检测或确认无信号还需 \(5\) 秒；这并不是完成定位与清除的总时间。

若任务要求第二次绝不丢失信号，使用约 \((864.3,\pm511.4)\) 的方案，并接受约 \(5.94\%\) 的平均面积增加。该约束表示对每个可能源位置及其后验允许的半径都有接收保证；它不是简单要求到某个代表源点的距离小于 \(1000\) 米。

若以第三问的清除速度为最终目的，还应把实际反馈后的移动至包围圆圆心、光学尝试和继续检测动作纳入序贯决策。\(J\) 最优只保证当前单步准则的数值最优，不能推出整个任务耗时最短。主方案的 \(P_{20}\) 只有约 \(20.66\%\)，多数反馈还不能给出 \(20\) 米的统一覆盖保证。

\subsection{已验证事项与未决限制}
此次计算验证了代码在指定基准和所抽查案例下的积分收敛、镜像对称、概率归一化、后验采样覆盖及独立几何一致性。搜索采用有界全域网格及局部细化，\textbf{没有使用连续全局分支定界，也没有得到全局最优的形式化误差证书}；因此坐标应称为已搜索范围内稳定的数值最优解，而不是严格证明唯一的连续全局极小点。

源位置面积均匀、角误差均匀、不同地点误差独立、接收半径均匀都是约定，其中空间相关误差尚未作敏感性分析。假如题目的实际环境统计规律不同，应替换对应先验或似然后重新求解。本次仅作第二问数学模型的自主数值计算，未使用第三、四问模拟器或正式测试机会。

本次原点基准的南北镜像源于目标圆域与已知示向度的对称性；对任意非原点第一检测点，不能只平移本结论，必须重新计算第一支持集与 \(1800\) 米目标圆域的交集。所有细节可回溯到模型快照、代码、输入网格和核验文件。
\end{document}
